\subsection{Symplectic Classification and Transvections}

Throughout this subsection, symplectic geometries are finite-dimensional over a fixed arbitrary field $\mathbb{F}$.

\begin{theorem} \label{thm:symplectic-geometries-hyperbolic-classification} \hlt{Classification of Symplectic Geometries}\\
Let $\mathcal{V}$ be a symplectic geometry, $R=\operatorname{rad}(\mathcal{V})$, $n=\dim(\mathcal{V})$, and $r=\operatorname{rank}(\mathcal{V})$.
\begin{enumerate}[label=(\roman*)]
\setlength{\itemsep}{0pt}
\item If $R=\{\mathbf{0}\}$, then $\mathcal{V}$ is hyperbolic and $n$ is even.
\item $\mathcal{V}=R\odot\mathcal{H}$ for a hyperbolic subspace $\mathcal{H}\leq\mathcal{V}$, and $R$ is totally singular.
\item $r\in2\mathbb{Z}_{\geq0}$, and if $r=2k$, then $\mathcal{V}$ admits a basis $\beta$ satisfying $M_{\beta}=X_{2k,n-2k}\defeq\operatorname{diag}(J_2^{\oplus k},0_{n-2k})$, where $J_2\defeq\left(\begin{smallmatrix}0&1\\-1&0\end{smallmatrix}\right)$.
\item For symplectic geometries $\mathcal{V},\mathcal{W}$, $\mathcal{V}\cong\mathcal{W}\Longleftrightarrow(\dim\mathcal{V},\operatorname{rank}(\mathcal{V}))=(\dim\mathcal{W},\operatorname{rank}(\mathcal{W}))$.
\end{enumerate}
\begin{proof}
\begin{enumerate}[label=(\roman*)]
\setlength{\itemsep}{0pt}
\item[$(i)$] The case $n=0$ is immediate. Assume $n>0$ and fix $\mathbf{x}\neq\mathbf{0}$. By alternation and Theorem \ref{thm:isotropic-vector-sits-in-hyperbolic-plane}, $\mathbf{x}$ lies in a hyperbolic plane $\mathcal{H}_1$.\\
By Theorem \ref{thm:metric-orthogonal-complement-dimension}(ii),(iv), $\mathcal{V}=\mathcal{H}_1\odot\mathcal{H}_1^{\perp}$ with $\mathcal{H}_1^{\perp}$ nonsingular. Thus $\mathcal{V}$ is hyperbolic by induction on $n$, so $n$ is even.
\item[$(ii)$] Fix $\mathcal{V}=R\odot\mathcal{S}$ with $\mathcal{S}$ nonsingular and $R$ totally singular by Theorem \ref{thm:radical-splitting-metric-space}. The subspace $\mathcal{S}$ is hyperbolic by (i), so take $\mathcal{H}=\mathcal{S}$.
\item[$(iii)$] In (ii), $\operatorname{rank}(\mathcal{V})=\dim(\mathcal{H})=2k$ and $\dim R=n-2k$ by the zero form on $R$ and nonsingularity of $\mathcal{H}$.\\
A hyperbolic basis of $\mathcal{H}$ followed by any basis of $R$ is a basis $\beta$ of $\mathcal{V}$ with $M_{\beta}=X_{2k,n-2k}$.
\item[$(iv)$] ($\Rightarrow$) An isometry preserves dimension and rank by its invertibility and form preservation.
\item[$(iv)$] ($\Leftarrow$) Then $\mathcal{V}\cong\mathcal{W}$ by (iii) and Proposition \ref{prop:metric-form-matrix-recovery-and-congruence}(i) for the normal-form bases.
\end{enumerate}
\hlt{Structural Commentary}: Split the radical, decompose the regular part into planes, and read off the normal form.\\
\hlt{Key Insight}: Hyperbolic planes are the unique nonsingular symplectic blocks.
\end{proof}
\end{theorem}

\begin{definition} \label{def:symplectic-transvection} \hlt{Symplectic Transvection}\\
Let $\mathcal{V}$ be a nonsingular symplectic geometry.\\
For $\mathbf{v}\in\mathcal{V}\setminus\{\mathbf{0}\}$ and $a\in\mathbb{F}$, the \hlt{symplectic transvection} with direction $\mathbf{v}$ and parameter $a$ is $\tau_{\mathbf{v},a}:\mathcal{V}\to\mathcal{V}$, $\mathbf{x}\mapsto\mathbf{x}+a\innerprod{\mathbf{x},\mathbf{v}}\mathbf{v}$.
\end{definition}

\begin{proposition} \label{prop:symplectic-transvection-properties} \hlt{Basic Properties of Symplectic Transvections}\\
Let $\mathcal{V}$ be a nonsingular symplectic geometry, $0\neq\mathbf{v}\in\mathcal{V}$, $a,b\in\mathbb{F}$, $c\in\mathbb{F}^{\times}$, and $\sigma\in\operatorname{Sp}(\mathcal{V})$. Then:
\begin{enumerate}[label=(\roman*)]
\setlength{\itemsep}{0pt}
\item $\tau_{\mathbf{v},a}\in\operatorname{Sp}(\mathcal{V})$;
\item $\tau_{\mathbf{v},a}=I_{\mathcal{V}}\Longleftrightarrow a=0$;
\item $\tau_{\mathbf{v},a}\tau_{\mathbf{v},b}=\tau_{\mathbf{v},a+b}$ and $\tau_{\mathbf{v},a}^{-1}=\tau_{\mathbf{v},-a}$;
\item $\sigma\tau_{\mathbf{v},a}\sigma^{-1}=\tau_{\sigma\mathbf{v},a}$;
\item $\tau_{c\mathbf{v},a}=\tau_{\mathbf{v},ac^2}$.
\end{enumerate}
\begin{proof}
\begin{enumerate}[label=(\roman*)]
\setlength{\itemsep}{0pt}
\item[$(i)$] Set $T=\tau_{\mathbf{v},a}$ and $\lambda(\mathbf{x})=\innerprod{\mathbf{x},\mathbf{v}}$. By bilinearity and alternation, $(\forall\,\mathbf{x},\mathbf{y}\in\mathcal{V})\,\innerprod{T\mathbf{x},T\mathbf{y}}=\innerprod{\mathbf{x},\mathbf{y}}+a\lambda(\mathbf{y})\lambda(\mathbf{x})-a\lambda(\mathbf{x})\lambda(\mathbf{y})=\innerprod{\mathbf{x},\mathbf{y}}$.\\
The map $T$ is invertible by $(\forall\,\mathbf{x}\in\mathcal{V})\,\tau_{\mathbf{v},-a}\tau_{\mathbf{v},a}(\mathbf{x})=\mathbf{x}=\tau_{\mathbf{v},a}\tau_{\mathbf{v},-a}(\mathbf{x})$. Hence $T\in\operatorname{Sp}(\mathcal{V})$.
\item[$(ii)$] ($\Rightarrow$) If $\tau_{\mathbf{v},a}=I_{\mathcal{V}}$ and $a\neq0$, fix $\mathbf{x}$ with $\innerprod{\mathbf{x},\mathbf{v}}\neq0$ by nonsingularity. Then $\tau_{\mathbf{v},a}\mathbf{x}-\mathbf{x}=a\innerprod{\mathbf{x},\mathbf{v}}\mathbf{v}\neq\mathbf{0}$, a contradiction.
\item[$(ii)$] ($\Leftarrow$) If $a=0$, then $\tau_{\mathbf{v},a}=I_{\mathcal{V}}$.
\item[$(iii)$] By alternation, $(\forall\,\mathbf{x}\in\mathcal{V})\,\tau_{\mathbf{v},a}\tau_{\mathbf{v},b}(\mathbf{x})=\mathbf{x}+(a+b)\innerprod{\mathbf{x},\mathbf{v}}\mathbf{v}=\tau_{\mathbf{v},a+b}(\mathbf{x})$. Hence $\tau_{\mathbf{v},a}\tau_{\mathbf{v},b}=\tau_{\mathbf{v},a+b}$, and $\tau_{\mathbf{v},a}^{-1}=\tau_{\mathbf{v},-a}$ by $b=-a$.
\item[$(iv)$] By form preservation, $(\forall\,\mathbf{x}\in\mathcal{V})\,\sigma\tau_{\mathbf{v},a}\sigma^{-1}(\mathbf{x})=\mathbf{x}+a\innerprod{\sigma^{-1}\mathbf{x},\mathbf{v}}\sigma\mathbf{v}=\mathbf{x}+a\innerprod{\mathbf{x},\sigma\mathbf{v}}\sigma\mathbf{v}=\tau_{\sigma\mathbf{v},a}(\mathbf{x})$.
\item[$(v)$] By bilinearity, $(\forall\,\mathbf{x}\in\mathcal{V})\,\tau_{c\mathbf{v},a}(\mathbf{x})=\mathbf{x}+ac^2\innerprod{\mathbf{x},\mathbf{v}}\mathbf{v}=\tau_{\mathbf{v},ac^2}(\mathbf{x})$.
\end{enumerate}
\hlt{Structural Commentary}: Item (i) verifies the form; items (ii)--(v) compute the parameter laws.\\
\hlt{Key Insight}: Alternation cancels every quadratic correction term.
\end{proof}
\end{proposition}

\begin{theorem} \label{thm:symplectic-group-generated-by-transvections} \hlt{Symplectic Transvections Generate}\\
If $\mathcal{V}$ is nonsingular symplectic, then $\operatorname{Sp}(\mathcal{V})=\left\langle\tau_{\mathbf{v},a}:\mathbf{v}\in\mathcal{V}\setminus\{\mathbf{0}\},\ a\in\mathbb{F}\right\rangle$.
\begin{proof}
\begin{enumerate}[label=(\roman*)]
\setlength{\itemsep}{0pt}
\item \hlt{Generated Subgroup}: Set $G\defeq\left\langle\tau_{\mathbf{v},a}:\mathbf{v}\in\mathcal{V}\setminus\{\mathbf{0}\},\ a\in\mathbb{F}\right\rangle$. Then $G\leq\operatorname{Sp}(\mathcal{V})$ by Proposition \ref{prop:symplectic-transvection-properties}(i),(iii).
\item \hlt{Vector Transitivity}: Fix $\mathbf{x},\mathbf{y}\in\mathcal{V}\setminus\{\mathbf{0}\}$. If $\mathbf{x}=\mathbf{y}$, then $I_{\mathcal{V}}\mathbf{x}=\mathbf{y}$. If $\innerprod{\mathbf{x},\mathbf{y}}\neq0$, then $\tau_{\mathbf{y}-\mathbf{x},\innerprod{\mathbf{x},\mathbf{y}}^{-1}}(\mathbf{x})=\mathbf{y}$.\\
Suppose that $\innerprod{\mathbf{x},\mathbf{y}}=0$. If $\mathbf{y}=\lambda\mathbf{x}$ for $\lambda\in\mathbb{F}^{\times}$, fix $\mathbf{z}$ with $\innerprod{\mathbf{x},\mathbf{z}}\neq0$ by nonsingularity; then $\innerprod{\mathbf{z},\mathbf{y}}=-\lambda\innerprod{\mathbf{x},\mathbf{z}}\neq0$.\\
If $\mathbf{x},\mathbf{y}$ are independent, fix $\mathbf{z}$ with $\innerprod{\mathbf{x},\mathbf{z}}=\innerprod{\mathbf{z},\mathbf{y}}=1$ by Proposition \ref{prop:complements-exist-finite-dimensions} for extending the functional $f(\mathbf{x})=1$, $f(\mathbf{y})=-1$ and Theorem \ref{thm:riesz-map-for-nonsingular-metric-spaces}. In either case, set $c=\innerprod{\mathbf{x},\mathbf{z}}$ and $d=\innerprod{\mathbf{z},\mathbf{y}}$. Then $c,d\neq0$, $\tau_{\mathbf{z}-\mathbf{x},c^{-1}}(\mathbf{x})=\mathbf{z}$, and $\tau_{\mathbf{y}-\mathbf{z},d^{-1}}(\mathbf{z})=\mathbf{y}$.\\
Thus $G$ is transitive on $\mathcal{V}\setminus\{\mathbf{0}\}$.
\item \hlt{Hyperbolic-Pair Transitivity}: Fix hyperbolic pairs $(\mathbf{u},\mathbf{v})$ and $(\mathbf{u}',\mathbf{v}')$. Fix $\gamma\in G$ with $\gamma\mathbf{u}=\mathbf{u}'$ by (ii), and set $\mathbf{w}=\gamma\mathbf{v}$.\\
Then $(\mathbf{u}',\mathbf{w})$ is hyperbolic. If $\mathbf{w}=\mathbf{v}'$, take $\delta=I_{\mathcal{V}}$. If $\innerprod{\mathbf{w},\mathbf{v}'}\neq0$, take $\delta=\tau_{\mathbf{v}'-\mathbf{w},\innerprod{\mathbf{w},\mathbf{v}'}^{-1}}$; then $\delta\mathbf{u}'=\mathbf{u}'$ and $\delta\mathbf{w}=\mathbf{v}'$.\\
Suppose finally that $\mathbf{w}\neq\mathbf{v}'$ and $\innerprod{\mathbf{w},\mathbf{v}'}=0$. The vectors $\mathbf{u}',\mathbf{w},\mathbf{v}'$ are independent by pairing a relation with $\mathbf{u}'$ and $\mathbf{w}$. Fix $\mathbf{z}$ with $\innerprod{\mathbf{u}',\mathbf{z}}=\innerprod{\mathbf{w},\mathbf{z}}=\innerprod{\mathbf{z},\mathbf{v}'}=1$ by Proposition \ref{prop:complements-exist-finite-dimensions} for extending the functional satisfying $(f(\mathbf{u}'),f(\mathbf{w}),f(\mathbf{v}'))=(1,1,-1)$ and Theorem \ref{thm:riesz-map-for-nonsingular-metric-spaces}. Then $\delta=\tau_{\mathbf{v}'-\mathbf{z},1}\tau_{\mathbf{z}-\mathbf{w},1}$ fixes $\mathbf{u}'$ and maps $\mathbf{w}$ to $\mathbf{v}'$. In all cases, $\delta\gamma\in G$ maps the first pair to the second.
\item \hlt{Dimension Induction}: Fix $\sigma\in\operatorname{Sp}(\mathcal{V})$ and induct on $n=\dim\mathcal{V}$. If $n=0$, then $\sigma=I_{\mathcal{V}}\in G$.\\
For $n>0$, fix a hyperbolic plane $\mathcal{H}=\operatorname{span}(\mathbf{u},\mathbf{v})$ by Theorem \ref{thm:symplectic-geometries-hyperbolic-classification}(i), and fix $\gamma\in G$ satisfying $\gamma\sigma|_{\mathcal{H}}=I_{\mathcal{H}}$ by (iii).\\
The restriction of $\gamma\sigma$ to $\mathcal{H}^{\perp}$ is a product of transvections by induction.\\
Each such transvection extends to $\mathcal{V}$ by fixing $\mathcal{H}$ pointwise, so $\gamma\sigma\in G$ and $\sigma\in G$.
\end{enumerate}
\hlt{Structural Commentary}: Establish vector and pair transitivity, then remove one hyperbolic plane inductively.\\
\hlt{Key Insight}: Transvections move hyperbolic frames while fixing their orthogonal complements.
\end{proof}
\end{theorem}

\begin{proposition} \label{prop:symplectic-transformations-have-determinant-one} \hlt{Symplectic Transformations Have Determinant One}\\
If $\mathcal{V}$ is nonsingular symplectic, then $(\forall\,\sigma\in\operatorname{Sp}(\mathcal{V}))\,\det(\sigma)=1$.
\begin{proof}
By Theorem \ref{thm:symplectic-group-generated-by-transvections}, it suffices to prove $\det(\tau_{\mathbf{v},a})=1$. The map $\tau_{\mathbf{v},a}$ fixes $\mathbb{F}\mathbf{v}$ and induces $I_{\mathcal{V}/\mathbb{F}\mathbf{v}}$ by alternation.\\
Its matrix in a basis beginning with $\mathbf{v}$ is therefore triangular with unit diagonal, so $\det(\tau_{\mathbf{v},a})=1$.\\
\hlt{Structural Commentary}: The generation theorem reduces the determinant to one triangular rank-one map.\\
\hlt{Key Insight}: Every transvection has unit diagonal.
\end{proof}
\end{proposition}
